\documentclass[11pt]{article}
\usepackage[margin=1in]{geometry}
\usepackage{amsmath,amssymb}
\usepackage{booktabs}
\usepackage{hyperref}

\title{Potassium regulates axon-oligodendrocyte signaling and metabolic coupling in white matter}
\author{}
\date{}

\begin{document}
\maketitle

\begin{abstract}
The integrity of myelinated axons relies on homeostatic support from oligodendrocytes (OLs). How OLs detect axonal spiking and rapidly control axon--OL metabolic coupling is largely unknown. Combining optic nerve electrophysiology and two-photon imaging, we show that both high-frequency axonal firing and extracellular potassium (K$^+$) elevations trigger a fast Ca$^{2+}$ response in OLs, facilitated by barium-sensitive inwardly rectifying K$^+$ channels. Using OL-specific Kir4.1 knockout mice (Kir4.1 cKO) we show that oligodendroglial Kir4.1 regulates axonal energy metabolism and long-term axonal integrity in addition to K$^+$ clearance. Before axonal damage, young Kir4.1 cKO mice show reduced GLUT1 and MCT1 in myelin. Young Kir4.1 cKO optic nerve axons show lower resting lactate, diminished activity-induced lactate surges, and reduced glucose uptake and consumption. OLs detect high-frequency activity through K$^+$ signaling and this signalling tunes axon--OL metabolic coupling.
\end{abstract}

\section{Introduction}
Oligodendrocytes produce and maintain myelin and support axonal integrity \citep{r1,r2,r3,r4,r5,r6,r7,r8}. They are highly glycolytic and provide metabolic support via monocarboxylate transporters \citep{r9,r10,r11}, with surface GLUT1 regulated by glutamatergic signalling \citep{r12}. OLs also provide antioxidant support \citep{r21} and K$^+$ buffering \citep{r22,r23}. We hypothesised that activity-driven K$^+$ elevations regulate axon--OL metabolic coupling. Using optic nerve electrophysiology and two-photon imaging \citep{r24,r25}, we characterised OL Ca$^{2+}$ dynamics and axonal metabolite fluxes in wild-type and OL-specific Kir4.1 knockout mice.

\section{Methods}
\subsection{Mice and sensors}
PLP-CreERT;RCL-GCaMP6s mice \citep{r26,r27,r28} were treated with tamoxifen at 6--8 weeks; 3--5 month old mice were used. Kir4.1$^{fl/fl}$;MOGiCre (cKO) and Kir4.1$^{fl/fl}$ (ctrl) mice were used per \citep{r22}. AAV-mediated intravitreal delivery was used to express GCaMP6s (in axons), ATeam1.03 (ATP) \citep{r47}, Laconic (lactate) \citep{r49}, or FLII12Pglu700$\mu\Delta$6 (FLIIP) \citep{r51} in retinal ganglion cells and optic nerve axons.

\subsection{Electrophysiology and imaging}
Compound action potentials (CAPs) were recorded from optic nerves. 30\,s stimulation trains at 10, 25, or 50\,Hz alternated with 0.4\,Hz pulses. Pharmacology: TTX 1\,$\mu$M, NBQX 50\,$\mu$M, 7-CKA 100\,$\mu$M, D-AP5 100\,$\mu$M, PPADS 50\,$\mu$M, Suramin 50\,$\mu$M, Ba$^{2+}$ 100\,$\mu$M, Cd$^{2+}$, Ni$^{2+}$, nifedipine, benidipine, ruthenium red, ouabain 500\,$\mu$M, KB-R7943 25\,$\mu$M, SEA0400 10\,$\mu$M, Bumetanide 50\,$\mu$M, cytochalasin B 20\,$\mu$M, iodoacetate 1\,mM, NaN$_3$ 5\,mM. Glucose deprivation (GD) and mitochondrial inhibition (MI) were used to probe ATP recovery.

\subsection{Proteomics and immunoblotting}
TMT-based LC-MS/MS was performed on optic nerve lysates from 2.5-month cKO ($n=4$) and ctrl ($n=4$) mice. Purified myelin was immunoblotted for Kir4.1, MCT1, GLUT1, PLP, CNP, MOG, ATP1a1, and ATP1a3.

\subsection{Electron microscopy}
Optic nerves from 3-month and 7--8-month cKO and ctrl mice were processed for EM; g-ratios and axon size distributions were quantified.

\section{Results}
\subsection{Axonal activity evokes OL Ca$^{2+}$}
30\,s trains evoked a bi-phasic Ca$^{2+}$ response in OL somas, larger at higher frequencies ($n=108$ OLs from $N=8$ nerves; 50 vs.\ 25\,Hz $p<0.0001$; 50 vs.\ 10\,Hz $p<0.0001$; 25 vs.\ 10\,Hz $p<0.0001$; one-way ANOVA with Tukey). TTX (1\,$\mu$M) abolished the response (93\%$\pm$9 reduction, $n=82$, $N=6$, $p<0.0001$, paired $t$-test). Extracellular Ca$^{2+}$ removal (+200\,$\mu$M EGTA) reduced the response 96\%$\pm$14 ($n=44$, $N=3$, $p<0.0001$).

Bath-applied K$^+$ (5, 10, 30\,mM) evoked Ca$^{2+}$ responses that scaled with concentration (5 vs.\ 10\,mM $p=0.0048$; 5 vs.\ 30\,mM $p<0.0001$; 10 vs.\ 30\,mM $p<0.0001$; $n=35$--57, $N=4$--5). TTX did not change the K$^+$-evoked response ($n=72$, $N=5$, $p=0.8144$). Ba$^{2+}$ (100\,$\mu$M) reduced the stimulation-induced response by 84\%$\pm$10 ($n=45$, $N=4$, $p<0.0001$) and the K$^+$-evoked response by 88\%$\pm$9 ($n=47$, $N=3$, $p<0.0001$). Reverse-mode NCX blocker KB-R7943 (25\,$\mu$M) reduced the stimulation-evoked response by 44\%$\pm$11 ($n=64$, $N=5$, $p=0.0002$) and the K$^+$-evoked response by 47\%$\pm$8 ($n=52$, $N=3$, $p<0.0001$). AMPA blockade (NBQX) had no effect; NMDA blockade reduced the response by $\sim$20\%. PPADS and Suramin reduced it by $\sim$20\%. Bumetanide had no effect.

\subsection{OL-specific Kir4.1 is critical for K$^+$ clearance and long-term axonal integrity}
Young (2--3 month) Kir4.1 cKO nerves had normal CAP peak latencies (peak 1 $p=0.7637$; peak 2 $p=0.9958$; peak 3 $p=0.9265$), stimulus-response relationships, g-ratios ($p=0.1584$), and axon size distributions. CAP peak amplitude recovery was markedly slower in cKO and in wildtype + Ba$^{2+}$ (ctrl vs.\ cKO $p<0.0001$; wt vs.\ wt+Ba$^{2+}$ $p<0.0001$; ctrl vs.\ wt $p=0.9970$; cKO vs.\ wt+Ba$^{2+}$ $p=0.7362$). At 7--8 months, cKO showed significant axonal injury features ($n=5$ mice, 10 images each, 286\,$\mu$m$^2$; $p=0.0169$) and unique giant axonal swellings; GFAP increased but IBA1 did not.

\subsection{Kir4.1 cKO reduces GLUT1 and MCT1 in myelin}
TMT proteomics showed reduced Kir4.1 (Kcnj10), and top 50 FDR-sorted proteins revealed reductions in vesicular transport and energy metabolism. GSEA/KEGG indicated decreases in transmembrane transporter activity and oxidative phosphorylation. Immunoblot of purified myelin: Kir4.1 reduced by 81\%$\pm$13 ($p=0.0038$), GLUT1 reduced by 44\%$\pm$8 ($p=0.0044$), MCT1 reduced by 50\%$\pm$13 ($p=0.0179$).

\subsection{Minor ATP changes and reduced lactate in cKO axons}
Basal axonal ATP levels were comparable ($n=6$, $p=0.61$). ATP decline rate during GD+MI: $p=0.44$. ATP recovery rate: $p=0.51$. Onset of recovery was faster in cKO ($p=0.03$). Recovery magnitude was lower in cKO ($p=0.0135$). During 50\,Hz, ATP decline rate $p=0.81$; initial recovery rate $p=0.54$; long-term recovery slightly less in cKO ($p=0.0498$).

Lactate sensor calibration in wildtype confirmed dynamic range. Stimulation-evoked lactate surges scaled with frequency (50 vs.\ 25\,Hz $p=0.0141$; 50 vs.\ 10\,Hz $p<0.0001$; 25 vs.\ 10\,Hz $p=0.0141$). Ba$^{2+}$ reduced the stimulation-induced lactate rise by 34\%$\pm$11 ($n=7$, $p=0.0249$). In cKO vs.\ ctrl: basal lactate lower ($n=16$ vs.\ 8; $p=0.0119$); 50\,Hz-evoked lactate surge lower ($p=0.0174$); lactate decline during GD not significantly different ($p=0.1318$).

\subsection{Reduced axonal glucose metabolism in cKO}
Basal axonal glucose comparable between genotypes ($n=7$--9, $p=0.2276$). Iodoacetate-induced glucose rise slower in cKO by 36\%$\pm$13 ($n=4$--7, $p=0.0235$). Cytochalasin B glucose decline: ctrl $2.6\%\pm0.4$/min vs.\ cKO $1.6\%\pm0.3$/min, basal glucose consumption reduced 37\%$\pm$17 in cKO ($p=0.0384$). During 50\,Hz, glucose declined at $2.7\%\pm0.6$/min in ctrl and was essentially flat ($0.1\%\pm0.2$/min) in cKO ($p=0.0002$). Post-stimulation glucose rose above baseline more in cKO (3.2\%$\pm$0.3) than ctrl (1.4\%$\pm$0.4, $p=0.0006$). Glycolytic activation at 50\,Hz: ctrl 0.1 vs.\ 50\,Hz $p=0.0022$; cKO 0.1 vs.\ 50\,Hz $p<0.0001$; 50\,Hz glucose consumption reduced 40\%$\pm$15 in cKO ($p=0.0208$); fold change 0.1 $\to$ 50\,Hz: $8.5\pm 2$ ctrl vs.\ $8.5\pm 1$ cKO ($p=0.9858$).

\begin{table}[h]
\centering
\caption{Key effects in young Kir4.1 cKO vs.\ ctrl (mean $\pm$ SEM).}
\begin{tabular}{lcc}
\toprule
Measure & Result & $p$-value \\
\midrule
Myelin Kir4.1 abundance reduction & 81\%$\pm$13 & 0.0038 \\
Myelin GLUT1 reduction & 44\%$\pm$8 & 0.0044 \\
Myelin MCT1 reduction & 50\%$\pm$13 & 0.0179 \\
Basal glucose consumption reduction & 37\%$\pm$17 & 0.0384 \\
50\,Hz glucose consumption reduction & 40\%$\pm$15 & 0.0208 \\
\bottomrule
\end{tabular}
\end{table}

\section{Discussion}
Our findings provide a working model in which axon--OL signaling and metabolic coupling in white matter is controlled by extracellular K$^+$ homeostasis via oligodendroglial Kir4.1. High-frequency axonal firing raises extracellular K$^+$, depolarising OLs via Ba$^{2+}$-sensitive Kir channels and driving Ca$^{2+}$ entry primarily through reverse-mode NCX activation. OL Kir4.1 is primarily responsible for K$^+$ clearance in adult white matter: OL-specific Kir4.1 deletion and Ba$^{2+}$ treatment of wildtype nerves produce equivalent CAP recovery deficits after high-frequency stimulation. GLUT1 and MCT1 abundance in myelin is reduced in young cKO mice, 5 months before overt axonopathy, paralleled by diminished axonal lactate and glucose metabolism. This Kir4.1-mediated axon--OL crosstalk regulates surface expression of MCT1 and GLUT1 in myelin, suggesting a feed-forward regulation of metabolite supply to axons. Our results reveal for the first time that OLs regulate axonal glucose metabolism, a novel interaction that requires Kir4.1 function and may impact long-term axonal integrity via effects on the pentose phosphate pathway, biosynthesis, and redox state.

\bibliographystyle{plain}
\bibliography{references}

\end{document}
